\documentclass[11pt]{article}
\usepackage[margin=1in]{geometry}
\usepackage{amsmath,amssymb,amsthm,booktabs,array}
\usepackage[T1]{fontenc}
\usepackage{lmodern}
\usepackage[colorlinks=true,urlcolor=blue,citecolor=blue,linkcolor=blue]{hyperref}
\newtheorem{theorem}{Theorem}
\newtheorem{proposition}[theorem]{Proposition}
\newtheorem{corollary}[theorem]{Corollary}
\newcommand{\ph}{\varphi}
\title{Adamson's Fibonacci divisor-transform conjecture:\\
sharp ratio errors extending Lang's formula}
\author{Haobo Ma\\\small ChronoAI Pte Ltd, Singapore
\and Wenlin Zhang\\\small National University of Singapore\\\small ChronoAI Pte Ltd, Singapore}
\date{November 14, 2026}
\begin{document}
\maketitle
\begin{abstract}
Extending Wolfdieter Lang's odd-Fibonacci ratio formula (OEIS A001519,
COMMENTS, March 30, 2020), we prove Gary W. Adamson's conjecture,
``$a(n)/a(n-1)$ tends to $\ph^2$'' (OEIS A130095, COMMENTS,
original revision \#1, May 11, 2007; retained in revision \#11,
December 11, 2019), where $a(n)=\sum_{d\mid n}F_{2d-1}$ and
$\ph=(1+\sqrt5)/2$.
We obtain the sharper limits
$\ph^{2m}(a(2m)/a(2m-1)-\ph^2)\to\ph^2$ and
$\ph^{2m+1}(a(2m+1)/a(2m)-\ph^2)\to-\ph^3$.
The proof separates the largest proper divisor from an exponentially
smaller remainder. We also prove Bala's generating functions and prime
congruences, and give an explicit proof of Kote\v{s}ovec's previously
recorded Fibonacci-transform asymptotic.
\end{abstract}

\section{The sequence and its ratio}
Let $F_0=0$, $F_1=1$, and $F_{j+2}=F_{j+1}+F_j$ for $j\geq0$, and put
$\ph=(1+\sqrt5)/2$. For $n\geq1$ define
\begin{equation}\label{eq:a}
 a(n)=\sum_{d\mid n} F_{2d-1},
\end{equation}
where divisors are positive. This is Bala's formula for A130095
\cite{adamson}; it begins $1,3,6,16,35,97$. Adamson's comment states
``Conjecture: a(n)/a(n-1) tends to phi\^{}2.'' The comment already occurs in
the entry's original revision \#1, dated May 11, 2007, whose author
field credits Adamson, May 6, 2007. The later divisor-sum and
generating-function formulas are signed by Bala, March 26, 2015.

The divisor transform in \eqref{eq:a} adds many smaller-index terms,
but no proper divisor of $n$ exceeds $n/2$. The following estimate
makes their total size explicit.

\begin{theorem}\label{thm:main}
For every $n\geq1$,
\begin{equation}\label{eq:error}
 0\leq a(n)-F_{2n-1}\leq F_{2\lfloor n/2\rfloor}\leq F_n,
 \qquad
 0\leq \frac{a(n)}{F_{2n-1}}-1\leq\frac1{F_n}.
\end{equation}
Consequently
\begin{equation}\label{eq:limit}
 a(n)\sim\frac{\ph^{2n-1}}{\sqrt5},\qquad
 \lim_{n\to\infty}\frac{a(n)}{a(n-1)}=\ph^2.
\end{equation}
The ratio is defined for every $n\geq2$. Equivalently, for every
$\epsilon>0$ there is an integer $N\geq2$ such that
$|a(n)/a(n-1)-\ph^2|<\epsilon$ for every integer $n\geq N$.
\end{theorem}

\begin{proof}
Write $a(n)=F_{2n-1}+R(n)$, with $R(n)$ the sum over proper positive
divisors. Put $m=\lfloor n/2\rfloor$. Since the summands are
nonnegative, telescoping gives
\[
 0\leq R(n)\leq\sum_{d=1}^{m}F_{2d-1}
 =\sum_{d=1}^{m}(F_{2d}-F_{2d-2})=F_{2m}\leq F_n.
\]
For $n=1$ this sum is empty and equals $F_0=0$.
The Fibonacci addition identity
$F_{r+s}=F_{r-1}F_s+F_rF_{s+1}$ for $r\geq1$, $s\geq0$
follows by induction on $s$ from the recurrence. Taking $r=n$,
$s=n-1$ gives
\[
 F_{2n-1}=F_n^2+F_{n-1}^2\geq F_n^2.
\]
As $F_n>0$ for $n\geq1$, division proves \eqref{eq:error}.

For completeness, put $\psi=(1-\sqrt5)/2=-\ph^{-1}$. The expression
$(\ph^j-\psi^j)/\sqrt5$ satisfies the Fibonacci recurrence and its
initial values, so
\begin{equation}\label{eq:binet}
 F_j=\frac{\ph^j-\psi^j}{\sqrt5}.
\end{equation}
In particular $F_n\to\infty$ and
$F_{2n-1}\sim\ph^{2n-1}/\sqrt5$. Let
$e_n=R(n)/F_{2n-1}$. Then $e_n\to0$ by \eqref{eq:error}, and for
$n\geq2$,
\[
 \frac{a(n)}{a(n-1)}
 =\frac{F_{2n-1}}{F_{2n-3}}\frac{1+e_n}{1+e_{n-1}}.
\]
All denominators are positive. Formula \eqref{eq:binet} shows that the
first factor tends to $\ph^2$, and the second tends to $1$. This proves
both conclusions in \eqref{eq:limit} and the stated quantified form.
\end{proof}

In fact \eqref{eq:binet} and \eqref{eq:error} give
$a(n)/F_{2n-1}=1+O(\ph^{-n})$ and
$a(n)/a(n-1)=\ph^2+O(\ph^{-n})$.
For the second estimate, the odd-Fibonacci ratio differs from
$\ph^2$ by $O(\ph^{-4n})$, and
\[
 \left|\frac{1+e_n}{1+e_{n-1}}-1\right|
 \leq e_n+e_{n-1}=O(\ph^{-n}).
\]
Separating the divisor $n/2$ gives exact limiting constants for both
errors. Write $\eta_n=1$ for even $n$ and $\eta_n=0$ for odd $n$.

\begin{corollary}\label{cor:sharp}
For $n\geq2$,
\[
0\leq R(n)-\eta_n F_{n-1}\leq F_{2\lfloor n/3\rfloor}.
\]
Consequently
\[
\frac{a(n)}{F_{2n-1}}-1
=\eta_n\ph^{-n}+O(\ph^{-4n/3}),
\]
and
\begin{align*}
\lim_{m\to\infty}\ph^{2m}
 \left(\frac{a(2m)}{a(2m-1)}-\ph^2\right)&=\ph^2,\\
\lim_{m\to\infty}\ph^{2m+1}
 \left(\frac{a(2m+1)}{a(2m)}-\ph^2\right)&=-\ph^3.
\end{align*}
Both $O(\ph^{-n})$ error orders are therefore sharp. The successive
ratios eventually lie above $\ph^2$ at even indices and below it at
odd indices.
\end{corollary}
\begin{proof}
Every proper divisor other than $n/2$ has complementary divisor at
least three, hence is at most $n/3$. Bound their sum by the sum of all
odd Fibonacci terms through that index, as in Theorem~\ref{thm:main}.
Binet's formula gives
$F_{n-1}/F_{2n-1}=\ph^{-n}(1+O(\ph^{-2n}))$ and
$F_{2\lfloor n/3\rfloor}/F_{2n-1}=O(\ph^{-4n/3})$.
Thus, writing $e_n=R(n)/F_{2n-1}$,
\[
\frac{1+e_n}{1+e_{n-1}}
=1+e_n-e_{n-1}+O(\ph^{-2n})
=1+\eta_n\ph^{-n}-\eta_{n-1}\ph^{-(n-1)}
 +O(\ph^{-4n/3}).
\]
Multiply by $F_{2n-1}/F_{2n-3}=\ph^2+O(\ph^{-4n})$ and
then by $\ph^n$. The two parities give the displayed limits.
The even normalized-term limit is $1$, by the preceding expansion.
\end{proof}

\section{A general estimate and related formulas}
\begin{proposition}\label{prop:general}
Suppose $b_n\geq0$ for $n\geq1$ and $b_n\sim C\lambda^n$, where
$C>0$ and $\lambda>1$. For $A(n)=\sum_{d\mid n}b_d$,
\[
 A(n)=b_n\bigl(1+O(\lambda^{-n/2})\bigr),\qquad
 A(n)\sim C\lambda^n,\qquad
 \frac{A(n)}{A(n-1)}\longrightarrow\lambda.
\]
The quotient is defined for all sufficiently large $n$.
\end{proposition}
\begin{proof}
Choose $K>0$ so that $b_d\leq K\lambda^d$ for every $d\geq1$;
the finitely many initial values can be absorbed in $K$.
For all sufficiently large $n$, $b_n\geq (C/2)\lambda^n>0$, while
\[
 0\leq A(n)-b_n\leq K\sum_{d=1}^{\lfloor n/2\rfloor}\lambda^d
 \leq\frac{K\lambda^{\lfloor n/2\rfloor+1}}{\lambda-1}.
\]
Dividing proves the relative bound. The asymptotic follows, and
$b_n/b_{n-1}\to\lambda$ follows directly from $b_n\sim C\lambda^n$.
The quotient of the two relative correction factors tends to $1$.
\end{proof}

\begin{corollary}\label{cor:standard}
For $c(n)=\sum_{d\mid n}F_d$, the sequence A007435,
\[
 c(n)\sim \frac{\ph^n}{\sqrt5},\qquad
 \frac{c(n)}{c(n-1)}\longrightarrow\ph.
\]
More generally, for fixed integers $k\geq1$ and $s$ with $k+s\geq1$,
the sequence $\sum_{d\mid n}F_{kd+s}$ is asymptotic to
$\ph^{kn+s}/\sqrt5$ and has successive-term ratio tending to $\ph^k$.
\end{corollary}
\begin{proof}
Apply Proposition~\ref{prop:general} and \eqref{eq:binet} with
$(C,\lambda)=(1/\sqrt5,\ph)$ and
$(\ph^s/\sqrt5,\ph^k)$, respectively. The index restriction makes
every summand defined and positive.
\end{proof}
The first asymptotic in Corollary~\ref{cor:standard} is recorded by
V\'aclav Kote\v{s}ovec in A007435, FORMULA, May 21, 2018
\cite{sloane}. Here it is a cited formula with an explicit proof,
not a new asymptotic claim.

The base-sequence ratio is also recorded with a proof calculation by
Wolfdieter Lang, March 30, 2020, in A001519, COMMENTS \cite{odd}.
Its matrix calculation has a sign misprint, which can be corrected
without changing the stated sequence or limit.
\begin{corollary}\label{cor:matrix}
Let $M=\left(\begin{smallmatrix}1&1\\1&2\end{smallmatrix}\right)$
and let $I$ be the $2\times2$ identity matrix. Then
$M^2=3M-I$, and, for $n\geq1$,
\[
(M^n)_{11}=F_{2n-1}
=\frac{\ph^{2n-1}+\ph^{1-2n}}{\sqrt5}.
\]
Consequently $(M^n)_{11}/(M^{n-1})_{11}\to\ph^2$.
\end{corollary}
\begin{proof}
Direct multiplication gives
$M^2=\left(\begin{smallmatrix}2&3\\3&5\end{smallmatrix}\right)=3M-I$,
rather than the source's displayed $3M+I$.
Thus $u_n=(M^n)_{11}$ satisfies $u_{n+2}=3u_{n+1}-u_n$,
with $u_1=1$, $u_2=2$, just as $F_{2n-1}$ does.
The displayed formula and limit follow from \eqref{eq:binet}.
In the same source comment, the prefactor written as $1/(2r+1)$,
where $r=\ph$, is $1/(2r-1)=1/\sqrt5$.
\end{proof}

\begin{corollary}[Bala's generating functions]\label{cor:gf}
With $B(x)=\sum_{n\geq1}F_{2n-1}x^n$, one has
$B(x)=x(1-x)/(1-3x+x^2)$ and both identities below.
\end{corollary}
\begin{proof}
The odd subsequence
(A001519 with positive indices \cite{odd}) obeys
$F_{2n+3}=3F_{2n+1}-F_{2n-1}$, by the Fibonacci recurrence. Its
first two terms give
\[
 B(x)=\frac{x(1-x)}{1-3x+x^2}.
\]
Rearranging coefficientwise finite sums proves Bala's Lambert-series
identity \cite{adamson}:
\begin{equation}\label{eq:gf}
 \sum_{n\geq1}a(n)x^n
 =\sum_{j\geq1}\frac{F_{2j-1}x^j}{1-x^j}
 =\sum_{j\geq1}\frac{x^j(1-x^j)}{1-3x^j+x^{2j}}.
\end{equation}
Indeed both expansions have coefficient
$\sum_{j\mid n}F_{2j-1}$ at $x^n$. Replacing $x$ by $x^2$ gives
Bala's second formula:
\[
 \sum_{n\geq1}a(n)x^{2n}
 =\sum_{j\geq1}
 \frac{x^j}{1/(x^j-x^{-j})-(x^j-x^{-j})}.
\]
Each summand on the right, after simplification, is
$x^{2j}(1-x^{2j})/(1-3x^{2j}+x^{4j})$.
These are identities of formal power series; they also hold as
absolutely convergent analytic series for $|x|<\ph^{-2}$ in
\eqref{eq:gf}, and $|x|<\ph^{-1}$ in the last display. To justify
convergence in \eqref{eq:gf}, use $F_{2j-1}=O(\ph^{2j})$ and
$1-|x|^j\geq1-|x|$; the other representations have the same
absolutely summable double series. Values at $x=0$ in the Laurent
form are understood through the simplified power series.
\end{proof}

\begin{corollary}[Related transform identities]\label{cor:related}
For $c(n)=\sum_{d\mid n}F_d$, Jovovic's ordinary generating function
and Gutkovskiy's logarithmic identity in A007435 are
\[
 \sum_{n\geq1}c(n)x^n
 =\sum_{j\geq1}\frac{x^j}{1-x^j-x^{2j}},\qquad
 \sum_{n\geq1}\frac{c(n)x^n}{n}
 =-\sum_{j\geq1}\frac{F_j}{j}\log(1-x^j).
\]
For $b_n=F_n$ or $b_n=F_{2n-1}$, and $A(n)=\sum_{d\mid n}b_d$,
the inverse formula is
\[
 b_n=\sum_{d\mid n}\mu(d)A(n/d).
\]
Here $\mu(d)=0$ if $d$ has a squared prime factor, and otherwise is
$(-1)^r$ if $d$ is a product of $r$ distinct primes.
\end{corollary}
\begin{proof}
The Fibonacci recurrence gives $\sum_{j\geq1}F_jx^j=x/(1-x-x^2)$.
Summing this series at $x,x^2,\ldots$ proves the first identity.
For the second, expand $-\log(1-y)=\sum_{m\geq1}y^m/m$;
the coefficient of $x^n$ on the right is
$\sum_{j\mid n}(F_j/j)/(n/j)=c(n)/n$.
Both are formal identities and converge absolutely for $|x|<\ph^{-1}$.
Finally, changing the order of the finite divisor sums gives the
coefficient of $b_e$ in the inverse formula as
$\sum_{d\mid n/e}\mu(d)$, equal to $1$ if $e=n$ and $0$ otherwise:
for an integer with $r>0$ distinct prime factors the sum is $(1-1)^r$.
This also proves the inverse of the divisibility matrix A051731.
\end{proof}

\begin{corollary}[Bala's and Somos's prime congruences]\label{cor:prime}
For a prime $p$, Bala's A130095 formula is
\[
a(p)\equiv
\begin{cases}
3& p\equiv2,3\pmod5,\\
2& p\equiv1,4\pmod5,\\
0& p=5
\end{cases}\pmod p.
\]
For $c(p)$, Somos's A007435 formula has respectively the residues
$0,2,1$ in these three cases.
\end{corollary}
\begin{proof}
First let $p$ be odd and different from $5$. In the algebra
$\mathbb F_p[t]/(t^2-5)$, write $s$ for the class of $t$ and set
$\alpha=(1+s)/2$, $\beta=(1-s)/2$. The elements $2$ and $s$
are invertible. The recurrence proves
$F_j=(\alpha^j-\beta^j)/s$ in this algebra.
Let $\chi=(5/p)\in\{1,-1\}$ be the Legendre symbol.
Euler's criterion and quadratic reciprocity give
$s^p=\chi s$ and $\chi=(p/5)$; see \cite[Theorem~4.1.7,
Corollary~4.2.3]{stein}.
The binomial theorem in characteristic $p$ now gives
$\alpha^p=(1+\chi s)/2$, $\beta^p=(1-\chi s)/2$.
It follows that
\[
 F_p\equiv\chi,\qquad
 F_{p+1}\equiv(1+\chi)/2,\qquad
 F_{p-1}\equiv(1-\chi)/2\pmod p.
\]
Using the addition identity from Theorem~\ref{thm:main},
\[
 a(p)=1+F_{2p-1}
 \equiv1+\chi^2+((1-\chi)/2)^2=(5-\chi)/2,
 \qquad c(p)\equiv1+\chi\pmod p.
\]
The squares modulo $5$ are $1$ and $4$, giving the asserted cases.
For $p=2$, $a(2)=3$ and $c(2)=2$; for $p=5$,
$a(5)=35$ and $c(5)=6$. These verify the exceptional primes directly.
\end{proof}

\section{Coverage and verification}
The source statements used here have the following scope. The
generating-function and standard-transform rows credit the formulas
already recorded in OEIS.
\begin{center}\small
\begin{tabular}{@{}p{0.44\linewidth}p{0.50\linewidth}@{}}
\toprule
Source item & Coverage \\
\midrule
A130095, COMMENTS: ratio conjecture &
Proved for every eventual integer index; Theorem~\ref{thm:main}. \\
A130095 ratio error (not posed in the source) &
Not posed; sharp parity limits proved in Corollary~\ref{cor:sharp}. \\
A130095, FORMULA: Bala's divisor sum &
Definition \eqref{eq:a}; the same positive-index convention. \\
A130095, FORMULA: Bala's two generating functions &
Reproved; prior formulas Bala, 2015; Corollary~\ref{cor:gf}. \\
A130095, FORMULA: Bala's prime congruences &
Proved; formula credited to Bala, 2015; Corollary~\ref{cor:prime}. \\
A007435, FORMULA: Kote\v{s}ovec's asymptotic &
Proved; prior formula Kote\v{s}ovec, May 21, 2018;
Corollary~\ref{cor:standard}. \\
A007435: generating functions and prime congruences &
Proved; prior formulas Jovovic, 2002; Somos, 2012;
Gutkovskiy, 2018; Corollaries~\ref{cor:related}, \ref{cor:prime}. \\
A001519: Lang's odd-Fibonacci ratio &
Reproved; prior proof Lang, 2020; Corollary~\ref{cor:matrix}. \\
A001519: Lang's displayed $M^2=3M+I$ and prefactor $1/(2r+1)$ &
Refuted-and-corrected: $M^2=3M-I$ and $1/(2r-1)$;
Corollary~\ref{cor:matrix}. \\
A051731: divisibility matrix \cite{divisibility} &
Action is \eqref{eq:a}; inverse proved in Corollary~\ref{cor:related}. \\
A001519 operator iterations, other combinatorial interpretations;
A007435 other sequence decompositions &
Not covered; distinct from the divisor estimates and identities here. \\
\bottomrule
\end{tabular}
\end{center}

The pinned Lean file
\href{https://github.com/the-omega-institute/trureturing/blob/d51b6958b717911d2705676a6e9bf473b3b6f324/D5/S3/Arith/Divisors/OeisA130095RatioLimit.lean}{\texttt{OeisA130095RatioLimit.lean}}
contains the definition \texttt{a} in the namespace
\texttt{D5.S3.Arith.Divisors.OeisA130095RatioLimit} and its theorem
\texttt{result}. That theorem states exactly the $\epsilon,N$
conclusion of Theorem~\ref{thm:main}, with real division and $N\geq2$.
The file's private lemmas \texttt{remainder\_le},
\texttt{fib\_square\_le}, and \texttt{full\_ratio\_limit} give the
remainder bound, square bound, and full ratio limit used in its proof.
The general transform result, sharp parity limits, matrix correction,
generating functions, inversion, and prime congruences have written
proofs; they are not assertions of the declaration \texttt{result}.
The accompanying \texttt{verify.py} checks finite integer identities,
the displayed inequalities, and the 30 published sequence values;
\texttt{checks.json} specifies its finite ranges. Those finite checks
are supplementary to the all-index proof.

\begin{thebibliography}{9}
\bibitem{adamson} G. W. Adamson,
\emph{A130095: Inverse M\"obius transform of odd-indexed Fibonacci numbers},
The On-Line Encyclopedia of Integer Sequences, revision \#11,
December 11, 2019.
\url{https://oeis.org/A130095}.
Original comment: \url{https://oeis.org/history?seq=A130095&start=10}.
Formulas by P. Bala, March 26, 2015.
\bibitem{sloane} N. J. A. Sloane,
\emph{A007435: Inverse Moebius transform of Fibonacci numbers},
The On-Line Encyclopedia of Integer Sequences, revision \#49,
April 18, 2025.
Asymptotic formula by V. Kote\v{s}ovec, May 21, 2018.
\url{https://oeis.org/A007435}.
\bibitem{odd} N. J. A. Sloane,
\emph{A001519: $a(n)=3a(n-1)-a(n-2)$ for $n\geq2$, with
$a(0)=a(1)=1$}, revision \#923,
October 26, 2026. \url{https://oeis.org/A001519}.
\bibitem{divisibility} K. Strassburger,
\emph{A051731: Triangle read by rows: $T(n,k)=1$ if $k$ divides $n$,
$T(n,k)=0$ otherwise, for $1\leq k\leq n$}, revision \#147,
October 4, 2026. \url{https://oeis.org/A051731}.
\bibitem{stein} W. Stein,
\emph{Elementary Number Theory: Primes, Congruences, and Secrets},
online text dated January 23, 2017, Theorem 4.1.7 and Corollary 4.2.3,
pp. 72--74. \url{https://wstein.org/ent/ent.pdf}.
\end{thebibliography}
\noindent\small Licensed under
\href{https://creativecommons.org/licenses/by/4.0/}{Creative Commons
Attribution 4.0 International (CC BY 4.0)}.
\end{document}
